\documentclass[11pt,a4paper]{ctexart}

\usepackage[a4paper,margin=0.82in]{geometry}
\usepackage{amsmath,amssymb,bm}
\usepackage{setspace}
\usepackage{fancyhdr}

\setstretch{1.05}
\setlength{\parindent}{0pt}
\setlength{\parskip}{0.35em}
\setlength{\abovedisplayskip}{7pt}
\setlength{\belowdisplayskip}{7pt}
\setlength{\abovedisplayshortskip}{5pt}
\setlength{\belowdisplayshortskip}{5pt}
\pagestyle{fancy}
\fancyhf{}
\rfoot{\thepage}
\renewcommand{\headrulewidth}{0pt}

\title{\vspace{-1.2cm}\textbf{CS310 自然语言处理}\\[3pt]\large Assignment 2 书面题答案\vspace{-0.4cm}}
\author{}
\date{}

\begin{document}
\maketitle
\thispagestyle{fancy}

\section*{1. 证明交叉熵损失满足 $J(y,\hat{y})=-\log \hat{y}_o$}

由于真实标签 $y$ 是关于正确 outside word $o$ 的 one-hot 向量，因此
\[
y_o = 1,\qquad y_w = 0 \quad (w \ne o).
\]

交叉熵定义为
\begin{align*}
J(y,\hat{y})
&= -\sum_{w} y_w \log \hat{y}_w \\
&= -\left(y_o \log \hat{y}_o + \sum_{w \ne o} y_w \log \hat{y}_w \right) \\
&= -\left(1 \cdot \log \hat{y}_o + \sum_{w \ne o} 0 \cdot \log \hat{y}_w \right) \\
&= -\log \hat{y}_o.
\end{align*}

因此，当目标分布是 one-hot 向量时，交叉熵只保留模型分配给正确 outside word 的对数概率。

\section*{2. 计算 $\frac{\partial J}{\partial v_c}$}

记
\[
z = U^\top v_c,\qquad \hat{y} = \operatorname{softmax}(z),\qquad
J(y,\hat{y}) = -\sum_w y_w \log \hat{y}_w.
\]

对于 softmax 加交叉熵这一组合，关于打分向量 $z$ 的梯度为
\[
\frac{\partial J}{\partial z} = \hat{y} - y.
\]

又因为
\[
z = U^\top v_c,
\]
根据链式法则可得
\begin{align*}
\frac{\partial J}{\partial v_c}
&= \frac{\partial z}{\partial v_c}\frac{\partial J}{\partial z} \\
&= U(\hat{y} - y).
\end{align*}

所以最终答案为
\[
\boxed{\frac{\partial J}{\partial v_c} = U(\hat{y} - y)}.
\]

若写成展开形式，也可以表示为
\[
\frac{\partial J}{\partial v_c} = \sum_w (\hat{y}_w - y_w)\,u_w.
\]

\section*{3. 计算每个 outside word 向量 $u_w$ 的梯度}

对任意 outside word 向量 $u_w$，其对应的分数为
\[
z_w = u_w^\top v_c.
\]

由链式法则，
\begin{align*}
\frac{\partial J}{\partial u_w}
&= \frac{\partial J}{\partial z_w}\frac{\partial z_w}{\partial u_w} \\
&= (\hat{y}_w - y_w)v_c.
\end{align*}

因此分两种情况：

\textbf{当 $w=o$ 时：}
\[
\frac{\partial J}{\partial u_o} = (\hat{y}_o - 1)v_c.
\]

\textbf{当 $w\ne o$ 时：}
\[
\frac{\partial J}{\partial u_w} = \hat{y}_w v_c.
\]

统一写法为
\[
\boxed{\frac{\partial J}{\partial u_w} = (\hat{y}_w - y_w)v_c}.
\]

\end{document}
